\appendix
\section{The pentagon relation} \label{appendix}

As we highlighted in Remark~\ref{rem:lambda}, under the use of super Ptolemy relations, the $\lambda$-lengths associated to arcs in a polygon are well-defined. To see this, it is sufficient to (1) observe that two consecutive applications of the super Ptolemy relation yields the original arc, and (2) show that the pentagon relation of five alternating applications of the super Ptolemy relation yields the original two arcs.

For (1), we consider the quadrilateral as in Figure~\ref{fig:super_ptolemy}, and flip the arc $e$.
Note that \[ef = ac + bd + \sqrt{abcd} \sigma \theta.\]
After this flip, if we apply the super Ptolemy relation to the diagonal $f$ in the resulting quadrilateral (and denote the resulting diagonal as $g$), we would have
\[fg = bd + ac + \sqrt{abcd} \sigma' \theta'.\]
However, using relation equation~\eqref{eq:mu_cross}, we can rewrite the latter as
\[fg = bd + ac + \sqrt{abcd} \sigma \theta,\] and solving for $g$, we obtain $g = e$ as desired.

The calculations to prove (2) are considerably more involved, and more easily shown using the notation of $\widetilde{\theta}_i$'s
that appeared in Corollary~\ref{cor:laurent}.

We will examine the sequence of flips illustrated in Figure~\ref{fig:pentagon_flips},
and apply the super Ptolemy relations one-by-one.
%
\begin {figure}[h!]
\centering
\begin {tikzpicture}[scale=1.09, every node/.style={scale=0.8}]
 \coordinate (v1) at (0,1);
 \coordinate (v2) at ({cos(pi/2 r + 1*2*pi/5 r)}, {sin(pi/2 r + 1*2*pi/5 r)});
 \coordinate (v3) at ({cos(pi/2 r + 2*2*pi/5 r)}, {sin(pi/2 r + 2*2*pi/5 r)});
 \coordinate (v4) at ({cos(pi/2 r + 3*2*pi/5 r)}, {sin(pi/2 r + 3*2*pi/5 r)});
 \coordinate (v5) at ({cos(pi/2 r + 4*2*pi/5 r)}, {sin(pi/2 r + 4*2*pi/5 r)});

 \draw (v1) -- (v2) -- (v3) -- (v4) -- (v5) -- cycle;

 \draw[->-] (v1) -- (v3);
 \draw[->-] (v1) -- (v4);

 \draw (-0.6,0.8) node {\small $a$};
 \draw (0.6,0.8) node {\small $b$};
 \draw (0.9,-0.3) node {\small $c$};
 \draw (0,-1) node {\small $d$};
 \draw (-0.9,-0.3) node {\small $e$};

 \draw (-0.4,0) node {\small $x_1$};
 \draw (0.4,0) node {\small $x_2$};

 \draw (-0.6,0.2) node {\small $\theta_1$};
 \draw (0,0) node {\small $\theta_2$};
 \draw (0.6,0.2) node {\small $\theta_3$};

 \draw (0,-1.5) node{starting};
 \draw (0,-1.8) node{configuration};

 %%%%%%%%%%%%%%%%%%%%%%%%%%%%%%%%%%%%%%%%%%%%%%%%%

 \begin {scope}[shift={(2.5,0)}]
 \coordinate (v1) at (0,1);
 \coordinate (v2) at ({cos(pi/2 r + 1*2*pi/5 r)}, {sin(pi/2 r + 1*2*pi/5 r)});
 \coordinate (v3) at ({cos(pi/2 r + 2*2*pi/5 r)}, {sin(pi/2 r + 2*2*pi/5 r)});
 \coordinate (v4) at ({cos(pi/2 r + 3*2*pi/5 r)}, {sin(pi/2 r + 3*2*pi/5 r)});
 \coordinate (v5) at ({cos(pi/2 r + 4*2*pi/5 r)}, {sin(pi/2 r + 4*2*pi/5 r)});

 \draw (v1) -- (v2) -- (v3) -- (v4) -- (v5) -- cycle;

 \draw[->-] (v2) -- (v4);
 \draw[->-] (v1) -- (v4);

 \draw (-0.6,0.8) node {\small $a$};
 \draw (0.6,0.8) node {\small $b$};
 \draw (0.9,-0.3) node {\small $c$};
 \draw[red] (0,-1) node {\small $d$};
 \draw (-0.9,-0.3) node {\small $e$};

 \draw (-0.4,0) node {\small $x_3$};
 \draw (0.4,0) node {\small $x_2$};

 \draw (-0.4,-0.4) node {\small $\theta_4$};
 \draw (0,0) node {\small $\theta_5$};
 \draw (0.6,0.2) node {\small $\theta_3$};

 \draw (0,-1.5) node{flip $x_1$};
 \end {scope}

 %%%%%%%%%%%%%%%%%%%%%%%%%%%%%%%%%%%%%%%%%%%%%%%%%

 \begin {scope}[shift={(5,0)}]
 \coordinate (v1) at (0,1);
 \coordinate (v2) at ({cos(pi/2 r + 1*2*pi/5 r)}, {sin(pi/2 r + 1*2*pi/5 r)});
 \coordinate (v3) at ({cos(pi/2 r + 2*2*pi/5 r)}, {sin(pi/2 r + 2*2*pi/5 r)});
 \coordinate (v4) at ({cos(pi/2 r + 3*2*pi/5 r)}, {sin(pi/2 r + 3*2*pi/5 r)});
 \coordinate (v5) at ({cos(pi/2 r + 4*2*pi/5 r)}, {sin(pi/2 r + 4*2*pi/5 r)});

 \draw (v1) -- (v2) -- (v3) -- (v4) -- (v5) -- cycle;

 \draw[->-] (v2) -- (v4);
 \draw[->-] (v2) -- (v5);

 \draw (-0.6,0.8) node {\small $a$};
 \draw (0.6,0.8) node {\small $b$};
 \draw[red] (0.9,-0.3) node {\small $c$};
 \draw[red] (0,-1) node {\small $d$};
 \draw (-0.9,-0.3) node {\small $e$};

 \draw (-0.4,0) node {\small $x_3$};
 \draw (-0.3,0.4) node {\small $x_4$};

 \draw (-0.4,-0.4) node {\small $\theta_4$};
 \draw (0,0) node {\small $\theta_7$};
 \draw (0,0.6) node {\small $\theta_6$};

 \draw (0,-1.5) node{flip $x_2$};
 \end {scope}

 %%%%%%%%%%%%%%%%%%%%%%%%%%%%%%%%%%%%%%%%%%%%%%%%%

 \begin {scope}[shift={(7.5,0)}]
 \coordinate (v1) at (0,1);
 \coordinate (v2) at ({cos(pi/2 r + 1*2*pi/5 r)}, {sin(pi/2 r + 1*2*pi/5 r)});
 \coordinate (v3) at ({cos(pi/2 r + 2*2*pi/5 r)}, {sin(pi/2 r + 2*2*pi/5 r)});
 \coordinate (v4) at ({cos(pi/2 r + 3*2*pi/5 r)}, {sin(pi/2 r + 3*2*pi/5 r)});
 \coordinate (v5) at ({cos(pi/2 r + 4*2*pi/5 r)}, {sin(pi/2 r + 4*2*pi/5 r)});

 \draw (v1) -- (v2) -- (v3) -- (v4) -- (v5) -- cycle;

 \draw[->-] (v3) -- (v5);
 \draw[->-] (v2) -- (v5);

 \draw (-0.6,0.8) node {\small $a$};
 \draw (0.6,0.8) node {\small $b$};
 \draw (0.9,-0.3) node {\small $c$};
 \draw[red] (0,-1) node {\small $d$};
 \draw (-0.9,-0.3) node {\small $e$};

 \draw (-0.1,-0.3) node {\small $x_5$};
 \draw (-0.3,0.4) node {\small $x_4$};

 \draw (0.4,-0.4) node {\small $\theta_8$};
 \draw (0,0) node {\small $\theta_9$};
 \draw (0,0.6) node {\small $\theta_6$};

 \draw (0,-1.5) node{flip $x_3$};
 \end {scope}

 %%%%%%%%%%%%%%%%%%%%%%%%%%%%%%%%%%%%%%%%%%%%%%%%%

 \begin {scope}[shift={(10,0)}]
 \coordinate (v1) at (0,1);
 \coordinate (v2) at ({cos(pi/2 r + 1*2*pi/5 r)}, {sin(pi/2 r + 1*2*pi/5 r)});
 \coordinate (v3) at ({cos(pi/2 r + 2*2*pi/5 r)}, {sin(pi/2 r + 2*2*pi/5 r)});
 \coordinate (v4) at ({cos(pi/2 r + 3*2*pi/5 r)}, {sin(pi/2 r + 3*2*pi/5 r)});
 \coordinate (v5) at ({cos(pi/2 r + 4*2*pi/5 r)}, {sin(pi/2 r + 4*2*pi/5 r)});

 \draw (v1) -- (v2) -- (v3) -- (v4) -- (v5) -- cycle;

 \draw[->-] (v3) -- (v5);
 \draw[->-] (v3) -- (v1);

 \draw (-0.6,0.8) node {\small $a$};
 \draw[red] (0.6,0.8) node {\small $b$};
 \draw (0.9,-0.3) node {\small $c$};
 \draw[red] (0,-1) node {\small $d$};
 \draw (-0.9,-0.3) node {\small $e$};

 \draw (-0.1,-0.3) node {\small $x_5$};
 \draw (-0.3,0.4) node {\small $x_6$};

 \draw (0.4,-0.4) node {\small $\theta_8$};
 \draw (0,0) node {\small $\theta_{11}$};
 \draw (-0.6,0.2) node {\small $\theta_{10}$};

 \draw (0,-1.5) node{flip $x_4$};
 \end {scope}

 %%%%%%
 \begin{scope}[shift={(12.5,0)}]
 	\coordinate (v1) at (0,1);
 \coordinate (v2) at ({cos(pi/2 r + 1*2*pi/5 r)}, {sin(pi/2 r + 1*2*pi/5 r)});
 \coordinate (v3) at ({cos(pi/2 r + 2*2*pi/5 r)}, {sin(pi/2 r + 2*2*pi/5 r)});
 \coordinate (v4) at ({cos(pi/2 r + 3*2*pi/5 r)}, {sin(pi/2 r + 3*2*pi/5 r)});
 \coordinate (v5) at ({cos(pi/2 r + 4*2*pi/5 r)}, {sin(pi/2 r + 4*2*pi/5 r)});

 \draw (v1) -- (v2) -- (v3) -- (v4) -- (v5) -- cycle;

 \draw[->-] (v3)--(v1);
 \draw[->-] (v4)--(v1);

 \draw (-0.6,0.8) node {\small $a$};
 \draw (0.6,0.8) node {\small $b$};
 \draw (0.9,-0.3) node {\small $c$};
 \draw (0,-1) node {\small ${\color{red} d}$};
 \draw (-0.9,-0.3) node {\small $e$};

 \draw (-0.3,0.4) node {\small $x_6$};
 \draw (0.3,-0.2) node {\small $x_7$};


 \draw (-0.6,0.2) node {\small $\theta_{10}$};
 \draw (0,0) node {\small $\theta_{12}$};
 \draw (0.6,0.2) node {\small $\theta_{13}$};

 \draw (0,-1.5) node{flip $x_5$};
 \end{scope}

\end{tikzpicture}
\caption{Flip sequence to verify the pentagon relation. Red labels indicate reversed orientation.}\label{fig:pentagon_flips}
\end{figure}

Define the ``\emph{modified}'' $\mu$-invariants (as in Corollary~\ref{cor:laurent})
\[ \til_1 = \sqrt{\frac{e}{ax_1}} \theta_1, \qquad \til_2 = \sqrt{\frac{d}{x_1x_2}} \theta_2, \qquad \til_3 =\sqrt{\frac{c}{bx_2}} \theta_3. \]

\textbf{First flip:} After flipping $x_1$, we get $x_3$:
\begin{align*}
 x_3 = \frac{ad + ex_2}{x_1} + \frac{\sqrt{adex_2}}{x_1} \theta_1 \theta_2
 = \frac{ad+ex_2}{x_1} + ax_2 \widetilde{\theta}_1 \widetilde{\theta}_2
\end{align*}
and the new $\theta$'s are
\begin{align*}
 \theta_4 = \frac{\sqrt{ad} \theta_1 - \sqrt{ex_2} \, \theta_2}{\sqrt{x_1x_3}}
 = \frac{1}{\sqrt{dex_3}} \big( ad \widetilde{\theta}_1 - ex_2 \widetilde{\theta}_2 \big)
\end{align*}
and
\begin{align*}
 \theta_5 = \frac{\sqrt{ad} \theta_2 + \sqrt{ex_2} \theta_1}{\sqrt{x_1x_3}}
 = \sqrt{\frac{ax_2}{x_3}} \big( \widetilde{\theta}_1 + \widetilde{\theta}_2 \big).
\end{align*}

\textbf{Second flip:} Flip $x_2$ to get $x_4$:
\begin{align*}
 x_4 = \frac{ac + bx_3}{x_2} + \frac{\sqrt{acbx_3}}{x_2} \theta_5 \theta_3
 = \frac{acx_1+abd+bex_2}{x_1x_2} + ab \big( \widetilde{\theta}_1 \widetilde{\theta}_2 + \widetilde{\theta}_1 \widetilde{\theta}_3 + \widetilde{\theta}_2 \widetilde{\theta}_3 \big).
\end{align*}
The new $\theta$'s are
\begin{align*}
 \theta_6 = \frac{\sqrt{ac} \theta_3 + \sqrt{bx_3} \theta_5}{\sqrt{x_2x_4}}
 = \sqrt{\frac{ab}{x_4}} \big( \widetilde{\theta}_1 + \widetilde{\theta}_2 + \widetilde{\theta}_3 \big)
\end{align*}
and
\begin{align*}
 \theta_7 = \frac{\sqrt{ac} \theta_5 - \sqrt{bx_3} \theta_3}{\sqrt{x_2x_4}}
 = \frac{1}{\sqrt{cx_3x_4}} \big( ac\big(\widetilde{\theta}_1 + \widetilde{\theta}_2\big) - bx_3 \widetilde{\theta}_3 \big).
\end{align*}

{\allowdisplaybreaks \textbf{Third flip:} Flip $x_3$ to get $x_5$:
\begin{align*}
 x_5={} & \frac{ce + dx_4}{x_3} + \frac{\sqrt{cdex_4}}{x_3} \, \theta_4 \theta_7 \\
 %
={} & \frac{cex_1 x_2 + dac x_1 + d a b d + dbex_2}{x_1 x_2 x_3} +
 \frac{dab}{x_3}( \widetilde{\theta}_1\widetilde{\theta}_2 + \widetilde{\theta}_1\widetilde{\theta}_3 +
 \widetilde{\theta}_2 \widetilde{\theta}_3) \\
 %
& {}+ { \frac{1}{x_3^2} (ad \widetilde{\theta}_1 - ex_2 \widetilde{\theta}_2)\left( ac(\widetilde{\theta}_1 + \widetilde{\theta}_2)
 - bx_3 \widetilde{\theta}_3\right) } \\
 %
={} & \left(\frac{ad + ex_2}{x_3}\right)\left(\frac{bd+cx_1}{x_1x_2}\right)
 + \left(\frac{dabx_3 +adac + ex_2 ac}{x_3^2}\right)\widetilde{\theta}_1 \widetilde{\theta}_2 \\
 & {}
 + \left(\frac{dab}{x_3} - \frac{adb}{x_3}\right) \widetilde{\theta}_1 \widetilde{\theta}_3
 + \left(\frac{dab + ex_2 b}{x_3}\right) \widetilde{\theta}_2 \widetilde{\theta}_3 \\
 %
={} &{ \left(\frac{ad + ex_2}{x_3}\right)\left(\frac{bd+cx_1}{x_1x_2}\right)
 + \left(\frac{dab + acx_1}{x_3}\right)\widetilde{\theta}_1 \widetilde{\theta}_2 %\hspace{5em}
 + %\hspace{8em}
 b x_1 \widetilde{\theta}_2 \widetilde{\theta}_3 }\footnotemark \\
={} & { \left(\frac{ad + ex_2 + a x_1x_2 \widetilde{\theta}_1 \widetilde{\theta}_2}{x_3}\right)\left(\frac{bd+cx_1}{x_1x_2}\right)
 %\hspace{8.95em}
 + %\hspace{8em}
 b x_1 \widetilde{\theta}_2 \widetilde{\theta}_3 } \\
={} & \frac{bd+cx_1}{x_2} + bx_1 \widetilde{\theta}_2 \widetilde{\theta}_3{}^{\ref{footnote1}}
\end{align*}
\footnotetext{\label{footnote1}Here we used $\big(ad + ex_2 + a x_1x_2 \widetilde{\theta}_1\widetilde{\theta}_2\big)/x_3 = x_1$.}
$\!\!$with new $\theta$'s
\begin{align*}
 \theta_8 = \frac{\sqrt{ce} \theta_4 - \sqrt{dx_4} b \theta_7}{\sqrt{x_3x_5}} b
 = \frac{1}{\sqrt{cdx_5}} \left( \frac{-c(ad+ex_2)}{x_3} b \widetilde{\theta}_2 + bd b \widetilde{\theta}_3 \right)
 = \frac{1}{\sqrt{cdx_5}} \big( bd \widetilde{\theta}_3 - cx_1 \widetilde{\theta}_2 \big)
\end{align*}
and
\begin{align*}
 \theta_9 &= \frac{\sqrt{ce} \theta_7 + \sqrt{dx_4} \theta_4}{\sqrt{x_3x_5}}
 = \frac{1}{x_3 \sqrt{ex_4x_5}} \big( a(ec+dx_4) {\widetilde{\theta}_1} + e(ac-x_2x_4) \widetilde{\theta}_2 - be x_3 \widetilde{\theta}_3 \big) \\
 &= {\frac{1}{x_3 \sqrt{ex_4x_5}} \big( \big(a x_3x_5 - aex_2b \til_2\til_3 \big) \widetilde{\theta}_1 - bex_3 \big(\widetilde{\theta}_2 + \widetilde{\theta}_3\big) + eabx_2 \widetilde{\theta}_1 \widetilde{\theta}_2 \widetilde{\theta}_3 \big) }\footnotemark \\
 &= {\frac{1}{\sqrt{ex_4x_5}} \big( a x_5 \widetilde{\theta}_1 - be \big(\widetilde{\theta}_2 + \widetilde{\theta}_3\big) \big). }{}^{\ref{footnote2}}
\end{align*}
\footnotetext{\label{footnote2}Here we used $x_2x_4 = ac + bx_3 + abx_2 \big(\widetilde{\theta}_1+\widetilde{\theta}_2\big)\widetilde{\theta}_3$ and
$x_3x_5 = ce + dx_4 + \big( ac (ad + ex_2)\widetilde{\theta}_1\widetilde{\theta}_2 - ad b x_3 \widetilde{\theta}_1\widetilde{\theta}_1 + ex_2bx_3 \widetilde{\theta}_1\widetilde{\theta}_3\big)/x_3$.}

\textbf{Fourth flip:} Finally, flip $x_4$ to get $x_6$:
\begin{align*}
 x_6 &= \frac{be + ax_5}{x_4} + \frac{\sqrt{abex_5}}{x_4} \theta_9 \theta_6
 \\&= \frac{bex_2 + abd + acx_1}{x_2 x_4}+\frac{abx_1}{x_4} \widetilde{\theta}_2 \widetilde{\theta}_3
 + \frac{ab}{x_4^2} \big( a x_5 \widetilde{\theta}_1 - be \big(\widetilde{\theta}_2 + \widetilde{\theta}_3\big)\big)
 \big(\widetilde{\theta}_1 + \widetilde{\theta}_2 + \widetilde{\theta}_3\big) \\
 &={ \frac{bex_2 + abd + acx_1}{x_2 x_4} + \frac{abx_1}{x_4} \widetilde{\theta}_2 \widetilde{\theta}_3
 + \frac{ab}{x_4} \left( \frac{a x_5 + be}{x_4}\right) \big(\widetilde{\theta}_1 \widetilde{\theta}_2 + \widetilde{\theta}_1 \widetilde{\theta}_3\big) } \\
 &= { x_1 - \frac{abx_1}{x_4} \big(\widetilde{\theta}_1 \widetilde{\theta}_2 + \widetilde{\theta}_1 \widetilde{\theta}_3 + \widetilde{\theta}_2 \widetilde{\theta}_3 \big) + \frac{abx_1}{x_4} \widetilde{\theta}_2 \widetilde{\theta}_3
 + \frac{ab}{x_4} \left( \frac{a x_5 + be}{x_4}\right)
 \big(\widetilde{\theta}_1 \widetilde{\theta}_2 + \widetilde{\theta}_1 \widetilde{\theta}_3\big)} \footnotemark \\
 &= { x_1 - \frac{abx_1}{x_4} \big(\widetilde{\theta}_1 \widetilde{\theta}_2 + \widetilde{\theta}_1 \widetilde{\theta}_3 + \widetilde{\theta}_2 \widetilde{\theta}_3 \big) + \frac{abx_1}{x_4} \widetilde{\theta}_2 \widetilde{\theta}_3} \\
 &\phantom{=} \ {} + { \frac{ab}{x_4} \left(x_1 - \frac{abx_1}{x_4} (\widetilde{\theta}_1 \widetilde{\theta}_2 + \widetilde{\theta}_1 \widetilde{\theta}_3 + \widetilde{\theta}_2 \widetilde{\theta}_3) + \frac{abx_1}{x_4} \widetilde{\theta}_2 \widetilde{\theta}_3
\right)
 \big(\widetilde{\theta}_1 \widetilde{\theta}_2 + \widetilde{\theta}_1 \widetilde{\theta}_3\big)} {}^{\ref{footnote3}} \\
 &= {x_1 - \frac{abx_1}{x_4} \big(\widetilde{\theta}_1 \widetilde{\theta}_2 + \widetilde{\theta}_1 \widetilde{\theta}_3 + \widetilde{\theta}_2 \widetilde{\theta}_3 \big) + \frac{abx_1}{x_4} \widetilde{\theta}_2 \widetilde{\theta}_3
 + \frac{abx_1 }{x_4} \big(\widetilde{\theta}_1 \widetilde{\theta}_2 + \widetilde{\theta}_1 \widetilde{\theta}_3\big)} \\
 &= x_1.
\end{align*}
\footnotetext{\label{footnote3}Here we used $x_1 = {(acx_1+abd+bex_2)}/{x_2x_4} + {ab x_1}\big( \widetilde{\theta}_1 \widetilde{\theta}_2 + \widetilde{\theta}_1 \widetilde{\theta}_3 + \widetilde{\theta}_2 \widetilde{\theta}_3 \big)/{x_4} $.}
$\!\!$The new $\theta$'s are
\begin{align*}
\theta_{10}&=\frac{\theta_6\sqrt{be}+\theta_9\sqrt{ax_5}}{\sqrt{x_4x_6}}\\
&=\sqrt{be\over x_4x_6}\cdot\sqrt{ab\over x_4}\big(\til_1+\til_2+\til_3\big)+\sqrt{ax_5\over x_4x_6}\cdot\sqrt{1\over{ex_4x_5}}\big(ax_5\til_1-be\til_2-be\til_3\big)\\
&=\sqrt{a\over ex_6}\left(be+ax_5\over x_4\right)\til_1\footnotemark
\\
&=\sqrt{a\over ex_6}x_6\til_1\\
&=\sqrt{ax_1\over e}\til_1\footnotemark\\
&=\theta_1
\end{align*}
and
\addtocounter{footnote}{-1}
\footnotetext{Here we used the non-obvious fact that $\theta_9\theta_6\theta_1=0$.}
\addtocounter{footnote}{+1}
\footnotetext{We use the equality $x_6=x_1$ here.}
\begin{align*}
\theta_{11}&=\frac{\theta_9\sqrt{be}-\theta_6\sqrt{ax_5}}{\sqrt{x_4x_6}}\\
&={\sqrt{b}\over x_4\sqrt{x_5x_6}}\big( a x_5 \widetilde{\theta}_1 - be \big(\widetilde{\theta}_2 + \widetilde{\theta}_3\big) \big)-{a\sqrt{bx_5}\over x_4\sqrt{x_6}}\big(\til_1+\til_2+\til_3\big)\\
&=-\sqrt{b\over x_5x_6}\left(be+ax_5\over x_4\right)\big(\til_2+\til_3\big)\footnotemark\\
&=-\sqrt{b x_6\over x_5}\big(\til_2+\til_3\big).
\end{align*}
\footnotetext{Similar to above, here we made use of $\theta_9\theta_6(\theta_2+\theta_3)=0$.}

\textbf{Fifth flip:} Finally, we flip $x_5$ to $x_7$ and get back to the original triangulation (with different orientation):
\begin{align*}
x_7&=	\frac{bd+cx_6}{x_5} + \frac{\sqrt{bcdx_6}}{x_5}\theta_8\theta_{11}\\
&=\frac{bd+cx_6}{x_5}+{\sqrt{bcdx_6}\over x_5}\cdot
\frac{1}{\sqrt{cdx_5}} \big( bd \widetilde{\theta}_3 - cx_1 \widetilde{\theta}_2 \big)\left(-\sqrt{b x_6\over x_5}\big(\til_2+\til_3\big)\right)\\
&=\frac{bd+cx_1}{x_5}+{bx_1\over x_5}\left(\frac{bd+cx_1}{x_5}\right)\til_2\til_3\\
&=\left(\frac{bd+cx_1}{x_5}\right)\left(1+{bx_1\over x_5}\til_2\til_3\right)\\
&=\left(x_2-\frac{bx_1x_2}{x_5}\til_2\til_3\right)\left(1+{bx_1\over x_5}\til_2\til_3\right)\footnotemark\\
&=x_2+\left(-{bx_1x_2\over x_5}+{bx_1x_2\over x_5}\right)\til_2\til_3\\
&=x_2.
\end{align*}
\footnotetext{\label{fn4}We use $x_2x_5=bd+cx_1+bx_1x_2x_5\til_2\til_3$ here.}
$\!\!$Finally, the new $\theta$'s are
\begin{align*}
\theta_{12}&={\theta_{11}\sqrt{bd}+\theta_8\sqrt{cx_1}\over {\sqrt{x_5x_2}}}\\
&=-\left({bd+cx_1}\over{x_5} \right)\sqrt{x_6\over x_7d}\ \til_2\\	
&=-\left(x_2-\frac{bx_1x_2}{x_5}\til_2\til_3\right)\sqrt{x_1\over x_2d}\ \til_2 {}^{\ref{fn4},\ref{fn5} }	\\
&=-\sqrt{x_2x_1\over d}\ \til_2\\
&=-\theta_2,
\\
\theta_{13}&={\theta_{8}\sqrt{bd}-\theta_{11}\sqrt{cx_1}\over {\sqrt{x_5x_2}}}\\
&=\frac{\sqrt{b}}{x_5\sqrt{cx_7}}\big(bd\til_3-cx_1\til_2\big)+\frac{x_1\sqrt{bc}}{x_5\sqrt{x_7}}\big(\til_2+\til_3\big)\\
&=\left(bd+cx_1\over x_5\right)\sqrt{b\over cx_2}\ \til_3 \footnotemark\\
&=\left(x_2-\frac{bx_1x_2}{x_5}\til_2\til_3\right)\sqrt{b\over cx_2}\ \til_3{}^{\ref{fn4} }\\
&=\sqrt{x_2b\over c}\ \til_3\\
&=\theta_3.
\end{align*}
\footnotetext{Here we use $x_7=x_2$.\label{fn5}}

We see in the end that $x_6 = x_1$, $x_7=x_2$. Note also that the edge $x_6$, $x_7$ are oriented opposite of $x_1$, $x_2$, and so the $\lambda_{ij}$'s
are independent of the orientations. We also see that we get the same $\mu$-invariants (up to sign), as $\theta_{10} = \theta_1$,
$\theta_{12} = -\theta_2$, and $\theta_{13} = \theta_3$.

}
